% Generated draft: requires mathematical and source review.
\documentclass[11pt]{article}
\usepackage[margin=25mm]{geometry}
\usepackage{fontspec}
\setmainfont{Latin Modern Roman}
\usepackage{amsmath,amssymb,mathtools,unicode-math}
\setmathfont{Latin Modern Math}
\usepackage{enumitem,needspace,xurl,hyperref}
\setlength{\parindent}{0pt}
\setlength{\parskip}{5pt}
\setlength{\emergencystretch}{3em}
\newcommand{\N}{\mathbb{N}}
\newcommand{\Z}{\mathbb{Z}}
\newcommand{\Q}{\mathbb{Q}}
\newcommand{\R}{\mathbb{R}}
\newcommand{\C}{\mathbb{C}}
\newcommand{\F}{\mathbb{F}}
\newcommand{\A}{\mathbb{A}}
\newcommand{\PP}{\mathbb{P}}
\newcommand{\E}{\mathbb{E}}
\newcommand{\eps}{\varepsilon}
\newcommand{\dd}{\,\mathrm{d}}
\begin{document}
\section*{UnsolvedMath ID: 10400040}
\textbf{Source problem code:} catalogue entry
\subsection{Definitions and mathematical statement}
Let \(\Z\mathcal K\) denote the free abelian group on the intended isotopy classes of knots. For a fixed local move \(m\), the standard local-move filtration \(F_d(\Z\mathcal K,m)\) is generated by alternating sums
\[
 \sum_{S\subseteq\{1,\ldots,d\}}(-1)^{|S|}K_S,
\]
where \(d\) applications of \(m\) have disjoint supports and \(K_S\) is obtained from the original knot by performing exactly the moves indexed by \(S\). In this convention \(F_0=\Z\mathcal K\) and \(F_{d+1}\subseteq F_d\). The excerpt does not include its filtration definition, so this standard convention needs source confirmation.

For every degree \(d\), calculate the abelian group
\[
 F_d(\Z\mathcal K,m)/F_{d+1}(\Z\mathcal K,m)
\]
for local moves such as:
\begin{enumerate}
\item the \(\#\) move;
\item the pass move;
\item the \(\Delta\) move;
\item the \(n\)-gon move, with \(n\) a move parameter distinct from the filtration degree \(d\).
\end{enumerate}
The \(\#\) symbol names a local move and is not replaced by the connected-sum operation. The diagrams and orientation/framing conventions for these moves, the allowed range of \(n\), and the knot category are not supplied. The coefficient group is \(\Z\): computing rational dimensions alone would not determine possible torsion in the requested quotient. The phrase ``such as'' leaves the programme open to other local moves, without changing the four explicit cases.
\subsection{Short English statement}
Compute the integer associated graded groups of the knot filtration generated by local moves, particularly the \(\#\), pass, \(\Delta\), and \(n\)-gon moves. The exact move diagrams and filtration convention must be recovered from the source.
\subsection{Sources}
\begin{itemize}
\item Local UnsolvedMath registry, permanent ID 10400040, source code catalogue entry; no proposer or proposal year supplied; update 2026-08-06.
\item Ohtsuki, \emph{Problems on invariants of knots and 3-manifolds} (2002), Problem 2.18, PDF page 41: \url{https://msp.org/gtm/2002/04/gtm-2002-04-024s.pdf}. Registry access date 2026-07-29; PDF font-structure extraction. Provenance distinguishes 231 numbered bold headings from surrounding references. Present primary-literature status and redistribution terms require review.
\item The supplied summary reports partial results for some moves but provides no exact references or graded-group calculations.
\end{itemize}
\end{document}
